\documentclass{article}
\usepackage{amsmath, amsthm, amssymb, booktabs}
\usepackage[margin=1.2in]{geometry}

\newtheorem{proposition}{Proposition}
\newtheorem{theorem}{Theorem}
\newtheorem{corollary}{Corollary}
\theoremstyle{definition}
\newtheorem{definition}{Definition}
\newtheorem{remark}{Remark}

\title{Empirical Calibration and Model Implications}
\date{}

\begin{document}
\maketitle

\noindent We calibrate the competition model using three years of Chatbot Arena
leaderboard data (May 2023 -- May 2026, 215 weekly snapshots), individual
Arena battle records (116{,}000 conversations, April--July 2025), training-compute
metadata for frontier models, and OpenRouter API token-usage data.
The notation follows the companion theory notes: $N$ firms, $K$ dimensions,
production function $f$, budget vector $B$, evaluator weights $w$, user
weights $v$.

\section{Win-Rate Decay of Frontier Models}

\subsection*{Empirical finding}

We track four landmark models from their debut through May 2026, measuring
their Elo advantage over the field median at each weekly Arena snapshot.
Define the \emph{Elo gap} $\Delta_i(t) = \text{Elo}_i(t) - \text{median}(t)$,
where the median is over all models present in that snapshot.

\begin{table}[h]
\centering
\begin{tabular}{lrrrrr}
\toprule
Model & Debut gap & Gap now & $\hat\sigma$ at debut & $\hat\sigma$ now
      & Linear slope \\
\midrule
GPT-4 (Mar.\ 2023)         & $+127$ & $-113$ & 0.58 & 0.43 & $-0.31$ Elo/day \\
Claude 3 Opus (Mar.\ 2024) & $+183$ & $ -57$ & 0.61 & 0.46 & $-0.32$ Elo/day \\
GPT-4o (May 2024)          & $+210$ & $ -19$ & 0.63 & 0.49 & $-0.32$ Elo/day \\
o1 (Dec.\ 2024)            & $+211$ & $+47$ & 0.63 & 0.53 & $-0.32$ Elo/day \\
\bottomrule
\end{tabular}
\caption{Win-rate ($\hat\sigma$) computed as $\sigma(\Delta_i/400)$.
The exponential time constant fitted to each gap series is $\tau \approx 330$
days (range 296--370); see Figure~\ref{fig:decay}.}
\label{tab:decay}
\end{table}

\noindent Two facts stand out.  First, the linear decay rate is essentially
\emph{identical} across all four models ($-0.32 \pm 0.01$ Elo/day,
or $-116$ Elo/year) even though their debut gaps differ by a factor of
$1.7\times$.  Second, the gap series is well-described by exponential
decay with $\tau \approx 1$ year, but is statistically indistinguishable
from a linear decay over the observed window.

\subsection*{Model implication: B grows exponentially}

Under the log production function $f(x) = a\log x$ and the proportional Nash
$x_{i,k}^* = w_k B_i$ (Theorem~1 in the theory notes), the Elo gap between a
model with budget $B_\mathrm{top}$ and a competitor with budget $B_c(t)$ is

\begin{equation}\label{eq:gap}
  \Delta(t) \;=\; a \log\!\frac{B_\mathrm{top}}{B_c(t)}
  \quad \text{(in output units; scale by 400 for Arena Elo)}.
\end{equation}

The gap decays linearly in $t$ if and only if $B_c(t)$ grows
\emph{exponentially}: $B_c(t) = B_c(0)\,e^{\alpha t}$, giving

\[
  \frac{d\Delta}{dt} = -a\,\alpha.
\]

The empirical rate $|d\Delta/dt| = 116$ Elo/year $= 0.29$ output-units/year
(using Elo scale 400) implies $a\,\alpha = 0.29$/yr.
If, consistent with industry estimates, compute doubles roughly every 9 months
($\alpha = \ln 2 / 0.75 \approx 0.92$/yr), then $a \approx 0.31$, meaning a
tenfold increase in training compute raises model quality by roughly $a \ln 10
\approx 0.72$ output units, or $\approx 290$ Elo points.

\begin{remark}[Scale invariance and the source of convergence]
Under $f = a\log$, win-rates between firm $i$ and firm $j$ equal
$\sigma\!\bigl(a\log(B_i/B_j)\bigr)$, which depends \emph{only on the ratio
$B_i/B_j$}, not on absolute budgets.  Industry-wide compute scaling that
raises all budgets proportionally does \emph{not} by itself change win-rates.
Convergence requires the budget ratio to narrow: new entrants are launched with
budgets closer to the current frontier, reducing $B_\mathrm{top}/B_c(t)$ over
time.  This is the mechanism consistent with the data: as compute becomes
cheaper, new models are trained with budgets that approach the frontier's
launch budget, eroding the incumbent's advantage.
\end{remark}

\begin{remark}[Universality of the decay rate]
Since $d\Delta/dt = -a\,\alpha$ depends only on the industry-level compute growth
rate $\alpha$ (not on which firm holds the lead), the model predicts that all
frontier models should decay at the same rate---consistent with the empirical
uniformity of $-0.32$ Elo/day across models that differ in provider, launch
date, and debut advantage.  The differing zero-crossing times arise solely from
the differing debut gaps: a model that launches with twice the advantage persists
for twice as long before reaching parity.
\end{remark}

\noindent Inverting~\eqref{eq:gap}, the implied budget ratios at each model's
debut are:

\begin{center}
\begin{tabular}{lcc}
\toprule
Model & Debut Elo gap & Implied $B_\mathrm{top}/B_\mathrm{median}$ \\
\midrule
GPT-4 (Mar.\ 2023) & 127 & $1.4\times$ \\
Claude 3 Opus (Mar.\ 2024) & 183 & $1.6\times$ \\
GPT-4o (May 2024) & 210 & $1.7\times$ \\
o1 (Dec.\ 2024)   & 211 & $1.7\times$ \\
\bottomrule
\end{tabular}
\end{center}

\noindent These ratios are plausible: frontier models are estimated to use
1.5--2$\times$ the training compute of the median model at their launch.


\section{The Production Function $f$}

\subsection*{Empirical finding}

Three independent pieces of evidence point to the log family.

\medskip\noindent\textbf{(A) Cross-sectional scaling.}
Fitting benchmark scores to $\log_{10}$ of parameter count across 27 non-fine-tuned
frontier models gives:
\[
  \text{MMLU} \approx 9.2 \cdot \log_{10}(N_B) + 56.6 \qquad (r^2 = 0.22),
\]
where $N_B$ is parameter count in billions.  A power-law fit $\text{MMLU}
\propto N_B^\alpha$ gives similar or lower $r^2$, and the sign is positive and
consistent with a concave, log-like response.

\medskip\noindent\textbf{(B) Frontier capability over time.}
Taking the best-performing model at each Arena snapshot, the frontier Elo
grows at a linear rate:
\[
  \text{Elo}_{\max}(t) \approx 141 \cdot t + c \qquad (r^2 = 0.95),
\]
rising from $\approx 1065$ (May 2023) to $\approx 1504$ (May 2026), a gain
of 439 Elo in three years.  A log-in-time fit gives $r^2 = 0.90$; linear
is marginally better over this window.

\medskip\noindent\textbf{(C) Consistency with power-law loss scaling.}
The canonical Chinchilla/Kaplan scaling law states that cross-entropy loss
decreases as a power law: $L \propto C^{-\alpha}$ with $\alpha \approx 0.05$.
Benchmark performance $P \in [0,1]$ relates to loss approximately as
$P \approx \Phi(-L/\sigma)$ for some threshold distribution.  In the regime
$P \in (0.2, 0.9)$ where the normal CDF is nearly linear in its argument,
$P \approx c_0 - c_1 \, L \approx c_0 + c_1 \alpha \log C$, i.e.,
\emph{performance is log-linear in compute}.  The log production function
is therefore the natural empirical form of the scaling law after
the loss-to-performance transformation.

\subsection*{Model implication: $f = a\log$ is well-specified}

If $B(t) \propto e^{\alpha t}$ (compute doubles every $\sim 9$ months), then
under $f(x) = a\log x$:
\[
  f(B(t)) = a\alpha t + \text{const},
\]
predicting frontier performance grows \emph{linearly in time}---exactly
what is observed ($r^2 = 0.95$).  A power-law production function
$f(x) = x^\beta$ would instead predict frontier performance growing as
$e^{\alpha\beta t}$ (exponential in time), which is not supported.

\begin{remark}
By Theorem~3 of the theory notes, $f = a\log$ is the \emph{unique} production
function giving directionally homogeneous Nash for all $(B, w)$.  The empirical
evidence that $f$ is approximately log therefore provides an independent
motivation for the homogeneity result: it is not merely a theoretical
convenience but reflects the functional form that fits AI training dynamics.
\end{remark}


\section{Optimal Evaluation Weights: $w$ vs.\ $v$}

\subsection*{Empirical finding}

We estimate the evaluator weight $w_k$ as the share of Arena category-leaderboard
votes on task $k$, and the user weight $v_k$ as the share of organic user
conversations in category $k$ from 116{,}000 individual Arena battles
(April--July 2025).  Conversations are labeled using the Arena category tags:
coding (29\%), math (6\%), creative writing (8\%), general instruction (57\%).

\begin{table}[h]
\centering
\begin{tabular}{lrrr}
\toprule
Category & $w_k$ (Arena LB) & $v_k$ (organic convs.) & $w_k - v_k$ \\
\midrule
Coding              & 0.21 & 0.29 & $-0.08$ \\
Math                & 0.15 & 0.06 & $+0.09$ \\
Creative writing    & 0.13 & 0.08 & $+0.05$ \\
General instruction & 0.51 & 0.57 & $-0.06$ \\
\bottomrule
\end{tabular}
\caption{Evaluator weights $w$ (Arena category leaderboard vote shares) vs.\
user weights $v$ (organic Arena conversation category labels, $n=116$k).
Arena over-weights math by $2.5\times$ and creative writing by $1.6\times$
relative to actual user demand; coding is under-weighted.}
\label{tab:w_v}
\end{table}

\subsection*{Model implication: optimal $w$ equals $v$, with a measurable welfare cost}

Under $f(x) = a\log x$ and the proportional Nash $x_{i,k}^* = w_k B_i$
(Theorem~1), the best performance on task $k$ is
\[
  \max_i f(x_{i,k}^*) = a\log(w_k B_{\max}) = a\log w_k + a\log B_{\max},
\]
where $B_{\max} = \max_i B_i$.  Social welfare is
\[
  \mathrm{SW}(w) = \sum_k v_k \bigl[a\log w_k + a\log B_{\max}\bigr]
  = a\log B_{\max} + a\sum_k v_k \log w_k.
\]
The second term is $-a \cdot \mathrm{KL}(v \,\|\, w) + \text{const}$, so
$\mathrm{SW}$ is maximized at $w^* = v$ (Gibbs inequality):

\begin{proposition}[Optimal evaluation weights]\label{prop:opt-w}
Under $f = a\log$ and the proportional Nash, social welfare is maximized over
all $w \in \Delta^{K-1}$ when $w_k = v_k$ for all $k$.  The welfare loss from
using $w \neq v$ is
\[
  \mathrm{SW}(v) - \mathrm{SW}(w) = a \cdot \mathrm{KL}(v \,\|\, w) \geq 0.
\]
\end{proposition}

\noindent That is, a benchmark whose category weights deviate from user demand
induces firms to over-invest in over-weighted tasks and under-invest in
under-weighted ones, reducing social welfare in proportion to the KL divergence.

\medskip
\noindent\textbf{Empirical magnitude.}
Using $w = w_\text{Arena}$ and $v = v_\text{conv}$ from Table~\ref{tab:w_v}:
\[
  \mathrm{KL}(v \,\|\, w_\text{Arena})
  = \sum_k v_k \log\frac{v_k}{w_k}
  \approx 0.062 \text{ nats},
\]
a welfare loss of approximately 6\% relative to the optimum.
The dominant contributions come from over-weighting math ($v_k = 0.06$,
$w_k = 0.15$) and under-weighting coding ($v_k = 0.29$, $w_k = 0.21$):
firms compete too hard on math benchmarks and not hard enough on coding,
relative to what users actually need.

\begin{remark}[OpenRouter versus Arena as proxies for $v$]
API usage data (OpenRouter, Jan 2025--Jun 2026) shows a different distribution:
$\approx$94\% of token volume flows through general-purpose models, with
$\approx$5\% through coding-specialized models (primarily grok-code).
This reflects that \emph{developers route coding tasks through general-purpose
models} rather than switching to a specialized one, so the model-level token
distribution collapses to ``general.''  The organic conversation labels
from Arena (Table~\ref{tab:w_v}) are a more accurate proxy for $v$ because
they measure what users \emph{ask for} rather than which model they choose to
ask.
\end{remark}

\begin{remark}[Richer social welfare]
The formula above assumes users select the best model \emph{per task} ($\max_i$),
which upper-bounds welfare.  If users commit to a single model for all tasks,
social welfare is $\sum_k v_k f(x_{i^*, k})$ where
$i^* = \arg\max_i \sum_k v_k f(x_{i,k})$.  In that case the firm with the
largest $B_i$ always wins, and the welfare analysis reduces to a question
about the right production function rather than evaluation weights.
The per-task-maximum formula is appropriate when users can switch models across
tasks, which is increasingly the default in agentic or multi-model pipelines.
\end{remark}


\section{Summary of Calibrated Parameters}

Table~\ref{tab:params} collects the empirically calibrated values for the
main model parameters.

\begin{table}[h]
\centering
\begin{tabular}{llll}
\toprule
Parameter & Symbol & Value & Source \\
\midrule
Production function & $f$ & $a\log x$ & Scaling law evidence (\S2) \\
Scale coefficient & $a$ & $\approx 0.31$ & Decay rate + compute doubling (\S1) \\
Compute growth rate & $\alpha$ & $\approx 0.92$/yr ($\times 2$ per 9 mo.) & Industry estimates \\
Elo decay rate & $|d\Delta/dt|$ & 116 Elo/yr ($= a\alpha \times 400$) & Arena time-series (\S1) \\
Frontier Elo growth & $\partial_t \text{Elo}_{\max}$ & $+141$/yr & Arena time-series (\S2) \\
Optimal evaluator weights & $w^*$ & $w^* = v$ & Proposition~\ref{prop:opt-w} \\
Arena--user KL divergence & $\mathrm{KL}(v\,\|\,w_\text{Arena})$ & 0.062 nats ($\approx 6\%$) & Table~\ref{tab:w_v} \\
\bottomrule
\end{tabular}
\caption{Calibrated parameters. The value $a \approx 0.31$ is derived from the
identity $a\alpha = 0.29$ (output-unit decay rate) and $\alpha = \ln 2/0.75$.}
\label{tab:params}
\end{table}

\end{document}
